\section{The Development Loop}
\label{sec:ch02-the-development-loop}

\index{schema!as a reviewed file}%
The loop that works is small. Change a type, export the schema, read the diff,
then send a request. The export is what makes it work, because it turns a
question you would otherwise answer by clicking around an IDE into a file your
version control system can answer for you. Commit \code{schema.graphql} and a
pull request that adds a field shows it as an added line, and a pull request
that removes one shows that too, in red, where a reviewer will see it.
Chapter~\ref{ch:breaking-changes} turns that from a habit into a gate.

The second half of the loop is asking. Point a browser at the same address the
requests go to and Hot Chocolate answers with Nitro, an IDE that reads the
schema and completes your query against it. It was called Banana Cake Pop
until October 2024, so anything written before then describes this same tool
under the older name. Rafael Staib announced the change on ChilliCream's blog,
writing that
\enquote{by renaming Banana Cake Pop and Barista to Nitro, we're simplifying
our ecosystem and making it easier for developers to navigate and interact
with our suite of products}~\autocite{staib2024nitro}.

\index{Nitro!serve mode}%
Where that IDE comes from is a default worth changing. Ask for it with the
stock configuration and the response carries somebody else's server header:

\begin{minted}{text}
HTTP/1.1 200 OK
Content-Type: text/html
Server: cloudflare
Last-Modified: Thu, 13 Aug 2026 10:13:03 GMT
\end{minted}

Your service is on localhost and cannot produce that header. It is fetching the
IDE from ChilliCream's content delivery network on request and passing it
through, which is what the default serve mode means: always the current build.
The package you installed also contains a copy, and serving that one instead is
one assignment, plus the two \code{using} directives that bring its two names
into scope.

\begin{minted}[highlightlines={1,2,12-19},fontsize=\footnotesize]{csharp}
using ChilliCream.Nitro.App;
using HotChocolate.AspNetCore;

var builder = WebApplication.CreateBuilder(args);

builder
    .AddGraphQL()
    .AddSessionsTypes();

var app = builder.Build();

app.MapGraphQL()
    .WithOptions(options =>
    {
        // Serve the Nitro build that shipped in the package. The
        // default, ServeMode.Latest, downloads the current one
        // from ChilliCream's CDN instead.
        options.Tool.ServeMode = ServeMode.Embedded;
    });

app.RunWithGraphQLCommands(args);
\end{minted}

Same request, same build, and now the service answers for itself:

\begin{minted}{text}
HTTP/1.1 200 OK
Content-Type: text/html
Server: Kestrel
Last-Modified: Tue, 09 Jun 2026 09:05:20 GMT
\end{minted}

\index{reproducibility}%
That second date is the assembly's, not today's. I pin this in every project I
start, so that the tool a colleague sees is the tool that came out of the
package we both restored rather than whatever shipped this morning. The
embedded copy is also the only one a disconnected machine could serve, though I
have not measured what the default does when the network is gone and this
chapter does not claim it fails.

\index{validation!before execution|(}%
The last part of the loop is the part that makes it fast, and it is easiest to
see by getting something wrong on purpose. Misspell a field:

\begin{minted}{text}
POST /graphql HTTP/1.1
Host: localhost:5001
Content-Type: application/json

{"query":"{ sessions { titel } }"}
\end{minted}

and back:

\begin{minted}[fontsize=\footnotesize, breakanywhere]{json}
{"errors":[{"message":"The field `titel` does not exist on the type `Session`.","locations":[{"line":1,"column":14}],"extensions":{"type":"Session","field":"titel","responseName":"titel","specifiedBy":"https://spec.graphql.org/September2025/#sec-Field-Selections"}}]}
\end{minted}

That answer arrives with status \code{400 Bad Request}, and it has no
\code{data} key at all, not even \code{null}. The status is the first surprise,
since a 400 from a GraphQL endpoint is easy to read as a transport failure when
it is a message about your document. The missing \code{data} key is the second,
and it is the one your client code has to handle.

What the response does not contain is any evidence that your code ran, because
none of it did. Figure~\ref{fig:ch02-two-requests} puts the two requests side
by side. The server parses the document, validates it against the schema built
at startup, and only then calls anything you wrote. A field name that is not in
the schema is settled before the first resolver, which is why the loop is quick:
most of what you get wrong is caught by a program that already knows the answer.

\begin{figure}
  \centering
  % Two requests against the same schema: one fails validation, one does not.
%
% Same idiom as figures/tikz/ch01-one-field.tex: each lane is a timeline read
% left to right, ticks mark stages, and a brace under the stages that matter
% carries a one-line label. Lane A and lane B are the same URL, the same
% service, the same HTTP method (POST /graphql). The parse and validate ticks
% sit at the same x in both lanes, so the reader sees identical machinery up
% to that point and only a divergence after validation runs.
%
% Bare tikzpicture (house style): the figure environment, caption and label
% live at the call site in the section file.
\begin{tikzpicture}[
  x=1mm, y=1mm,
  every node/.append style={font=\sffamily\scriptsize},
  axis/.style={draw, line width=0.4pt, -{Straight Barb[length=1.4mm, width=1.6mm]}},
  tick/.style={draw, line width=0.4pt},
  event/.style={anchor=south, align=center, text width=22mm, inner sep=1pt},
  span/.style={draw, line width=0.4pt},
  spanlabel/.style={anchor=north, align=center, inner sep=2pt},
  lane/.style={anchor=west, font=\sffamily\scriptsize\itshape},
]

% ------------------------------------------------------------------- lane A
\node[lane] at (0,13) {A. \{ sessions \{ titel \} \}, a field the schema does not have};

\draw[axis] (0,0) -- (150,0);
\foreach \x in {15,50,85,120}{\draw[tick] (\x,-1.5) -- (\x,1.5);}

\node[event] at (15,3)  {parses the\\document};
\node[event] at (50,3)  {validates against\\the schema};
\node[event, text width=26mm] at (85,3)  {validation fails:\\no field named titel};
\node[event, text width=26mm] at (120,3) {400 Bad Request:\\errors key, no data};

\draw[span] (62,-4) -- (62,-7) -- (133,-7) -- (133,-4);
\node[spanlabel] at (97,-8) {no resolver ran};

% ------------------------------------------------------------------- lane B
\begin{scope}[shift={(0,-46)}]
  \node[lane] at (0,13) {B. \{ sessions \{ id title startsAt speaker \{ name bio \} \} \}, fields it has};

  \draw[axis] (0,0) -- (150,0);
  \foreach \x in {15,50,74,96,118,140}{\draw[tick] (\x,-1.5) -- (\x,1.5);}

  \node[event] at (15,3)  {parses the\\document};
  \node[event] at (50,3)  {validates against\\the schema};
  \node[event, text width=16mm] at (74,3)  {validation\\passes};
  \node[event, text width=20mm] at (96,3)  {sessions resolver\\runs: 4 sessions};
  \node[event, text width=16mm] at (118,3) {speaker resolver\\runs, per session};
  \node[event, text width=16mm] at (140,3) {200 OK:\\data key present};

  \draw[span] (86,-4) -- (86,-7) -- (126,-7) -- (126,-4);
  \node[spanlabel] at (106,-8) {this is where your C\# runs};
\end{scope}

\end{tikzpicture}

  \caption{The same endpoint answering two requests. Everything up to
  validation is identical, and the lanes diverge only after it. A misspelled
  field is settled against the schema, so no resolver of yours is reached and
  there is nothing to debug in your own code. That boundary is why an
  implementation-first schema is worth reading before you go looking for a bug.}
  \label{fig:ch02-two-requests}
\end{figure}

The \code{specifiedBy} member of the error names the edition of the GraphQL
specification the server validated against. It is a useful thing to check when
a client and a server disagree about whether a document is legal, because they
may be reading different editions.
\index{validation!before execution|)}%

\begin{onhc14}
Every file in this chapter compiles unchanged on Hot Chocolate 14 except
\code{Program.cs}, and there the difference is one statement. On 14,
\code{WithOptions} takes the options object rather than an action, and the
serve mode enumeration lives in \code{HotChocolate.AspNetCore} rather than in
the Nitro package, so that \code{using} goes away. The whole file, on 14, with
what moved marked:

\begin{minted}[highlightlines={1,11-15}]{csharp}
using HotChocolate.AspNetCore;

var builder = WebApplication.CreateBuilder(args);

builder
    .AddGraphQL()
    .AddSessionsTypes();

var app = builder.Build();

app.MapGraphQL()
    .WithOptions(new GraphQLServerOptions
    {
        Tool = { ServeMode = GraphQLToolServeMode.Embedded }
    });

app.RunWithGraphQLCommands(args);
\end{minted}

The project targets \code{net8.0}, because 14 does not target .NET 10.
Upgrading in the other direction is easier than that implies: 16 still targets
\code{net8.0}, so moving off 14 does not oblige you to move framework at the
same time.

The schema these files produce on 14 is identical for \code{Query},
\code{Session} and \code{Speaker}, down to the order of the fields, and
\code{@cost} is on them there too. The exported file differs in the parts
nobody writes. The \code{DateTime} scalar is described differently and points
at a different specification URL, the definition of \code{@specifiedBy} is
written out where 16 omits it, and the directive definitions are printed on one
line each rather than broken across several. The file also begins with a UTF-8
byte order mark, the bytes \code{EF BB BF}, which 16 does not write; if a tool
downstream of you reads that file, it is worth knowing which of the two
produced it. Appendix~\ref{app:hc14} carries the full mapping.
\end{onhc14}
